\documentclass[11pt,a4paper]{article}

% --- 核心宏包与中文支持 ---
\usepackage[utf8]{inputenc}
\usepackage{xeCJK}
\usepackage{amsmath, amssymb, amsfonts}
\usepackage{listings}
\usepackage{xcolor}
\usepackage{geometry}
\usepackage{tcolorbox}
\usepackage{booktabs}
\usepackage{tabularx}
\usepackage{setspace}
\usepackage{enumitem}
\usepackage[hidelinks]{hyperref}
\usepackage{bookmark}
\usepackage{fancyhdr}

% --- PDF 书签与元数据 ---
\hypersetup{
    colorlinks=false,
    pdfborder={0 0 0},
    bookmarks=true,
    bookmarksnumbered=true,
    bookmarksopen=true,
    bookmarksopenlevel=2,
    pdfcreator={XeLaTeX},
    pdftitle={03 - Data Structures},
    pdfauthor={黄子扬},
}

% --- 页面美化设置 ---
\geometry{top=0.7in,bottom=0.8in,left=0.6in,right=0.6in}
\onehalfspacing
\tcbuselibrary{skins,breakable}
\setlength{\parindent}{0pt}
\setlength{\parskip}{0.6em}

% --- 代码样式配置 ---
\definecolor{mainblue}{RGB}{0,51,102}
\definecolor{examred}{RGB}{180,0,0}
\definecolor{codebg}{RGB}{248,248,248}
\definecolor{commentgreen}{RGB}{0,128,0}

\lstset{
    backgroundcolor=\color{codebg},
    basicstyle=\ttfamily\footnotesize,
    keywordstyle=\color{blue}\bfseries,
    commentstyle=\color{commentgreen},
    stringstyle=\color{orange},
    breaklines=true,
    showstringspaces=false,
    frame=leftline,
    xleftmargin=2em,
    numbers=left,
    numbersep=5pt,
    numberstyle=\tiny\color{gray},
    language=Python,
    morekeywords={None, True, False, in, del, for, if, else, def, return, sum, min, max, len, sorted, list, tuple, dict, set}
}

% --- 自定义教学框 ---
\newtcolorbox{concept}[1]{colback=blue!5, colframe=mainblue, fonttitle=\bfseries, title={{【Concept / 核心概念】#1}}, arc=3pt, breakable}
\newtcolorbox{warning}[1]{colback=red!5, colframe=examred, fonttitle=\bfseries, title={{【Warning / 避坑指南】#1}}, arc=3pt, breakable}
\newtcolorbox{examplebox}[1]{colback=green!2, colframe=green!40!black, fonttitle=\bfseries, title={{【Example / 实操案例】#1}}, arc=2pt, breakable}

% --- 页眉页脚 ---
\pagestyle{fancy}
\fancyhf{}
\fancyhead[L]{\small\bfseries 03 - Data Structures / 数据结构}
\fancyhead[R]{\small\thepage}
\renewcommand{\headrulewidth}{0.5pt}

% --- 文档信息 ---
\title{\Huge \bfseries 03 - Data Structures}
\author{\Large \textbf{哈尔滨工业大学 \quad 黄子扬}}
\date{2026年6月8日}

\begin{document}

\maketitle
\tableofcontents
\newpage

% ======================================================================
\section{Introduction / 什么是数据结构}
% ======================================================================

\begin{concept}{Data Structure / 数据结构}
    A data structure models a collection of data, such as a list of numbers, a row in a spreadsheet, or a record in a database. \\
    数据结构是对数据集合的建模。

    So far, you have been working with fundamental data types like \texttt{str}, \texttt{int}, and \texttt{float}. Many real-world problems are easier to solve when simple data types are combined into more complex data structures. \\
    Python has three built-in data structures that are the focus of this chapter: \textbf{Tuples}, \textbf{Lists}, and \textbf{Dictionaries}. (Plus \textbf{Sets} as a bonus.)
\end{concept}

\subsection{Recap: print() Advanced / print() 进阶参数}

\begin{concept}{sep and end Parameters / sep 和 end 参数}
    \texttt{print()} can print one or multiple objects. By default:
    \begin{itemize}
        \item Objects are separated by a whitespace (\texttt{sep=" "}).
        \item Output ends with a line break (\texttt{end="\textbackslash n"}).
    \end{itemize}
    These behaviors can be changed via the \texttt{sep} and \texttt{end} parameters.
\end{concept}

\begin{examplebox}{print() with sep and end / print() 的 sep 和 end 参数}
\begin{lstlisting}
# Default behavior
print("hello", "there")
# Output: hello there (separated by space, ends with newline)

# Custom sep and end
print("hello", "there", sep="", end="!")
# Output: hellothere! (no separator, ends with "!" instead of newline)
\end{lstlisting}
\end{examplebox}

% ======================================================================
\section{Tuples / 元组 (Parentheses: \texttt{()})}
% ======================================================================

\subsection{What is a Tuple? / 什么是元组？}

\begin{concept}{Sequence Type / 序列类型}
    A \textbf{sequence} is an ordered list of values. Each element in a sequence is assigned an integer, called an \textbf{index}, that determines the order in which the values appear. Just like strings, the index of the first value in a sequence is \texttt{0}. \\
    序列是有序的值列表。每个元素都有一个整数索引，从 0 开始。

    \textbf{Tuples are ordered} because their elements appear in a fixed order. The first element of \texttt{(1, 2, 3)} is 1, the second is 2, the third is 3.
\end{concept}

\subsection{Core Properties / 核心属性}

\begin{itemize}[itemsep=0.5em]
    \item \textbf{Ordered / 有序}：Elements appear in a fixed sequence.
    \item \textbf{Heterogeneous / 异构}：Can contain different types (e.g., \texttt{(1, "John", "Red")}).
    \item \textbf{Index starts at 0}：索引从 0 开始。
    \item \textbf{Immutable / 不可变}：Once created, you cannot change, add, or remove elements.
\end{itemize}

\begin{examplebox}{Tuple Basics / 元组基础}
\begin{lstlisting}
tuple1 = (1, 2, 3)
print(type(tuple1))             # <class 'tuple'>

# Heterogeneous / 异构：可以包含不同类型
tuple2 = (1, "John", "Red")
print(tuple2)                   # (1, 'John', 'Red')

# Empty tuple / 空元组
tuple3 = ()
print(len(tuple3))              # 0
\end{lstlisting}
\end{examplebox}

\subsection{Tuple Operations / 元组操作}

\begin{examplebox}{Length, Indexing, Slicing / 长度、索引、切片}
\begin{lstlisting}
tuple2 = (1, "John", "Red")

# Length / 长度
print(len(tuple2))              # 3

# Indexing / 索引（和字符串一样！）
print(tuple2[1])                # 'John'

# Slicing / 切片
print(tuple2[1:3])              # ('John', 'Red')
\end{lstlisting}
\end{examplebox}

\subsection{Checking Existence / 成员检查}

\begin{examplebox}{The in Keyword / in 关键字}
\begin{lstlisting}
tuple2 = (1, "John", "Red")

if "Red" in tuple2:
    print('There is red')       # This runs
else:
    print('There is no red')
\end{lstlisting}
\end{examplebox}

\subsection{Iteration / 元组遍历}

\begin{examplebox}{Tuples Are Iterable / 元组可迭代}
\begin{lstlisting}
tuple2 = (1, "John", "Red")
for item in tuple2:
    print(item)
# Output:
# 1
# John
# Red
\end{lstlisting}
\end{examplebox}

\begin{warning}{Immutability / 不可变性的后果}
    Like strings, tuples are \textbf{immutable}. You cannot change the value of an element once created. \\
    元组一旦创建就无法修改。试图执行 \texttt{tuple2[1] = "David"} 会抛出 \texttt{TypeError}: \texttt{'tuple' object does not support item assignment}.
\end{warning}

% ======================================================================
\section{Lists / 列表 (Square Brackets: \texttt{[]})}
% ======================================================================

\subsection{What is a List? / 什么是列表？}

\begin{concept}{Another Sequence Type / 另一个序列类型}
    The list data structure is another sequence type in Python. Just like strings and tuples, lists contain items that are indexed by integers, starting with 0.

    A list literal looks almost exactly like a tuple literal, except it uses \textbf{square brackets} \texttt{[ ]} instead of parentheses.
\end{concept}

\begin{examplebox}{Creating Lists / 创建列表}
\begin{lstlisting}
mylist1 = [1, 2, 3]
print(type(mylist1))            # <class 'list'>

# Like tuples, lists can contain different types
mylist2 = [1, "John", "Red"]
print(mylist2)                  # [1, 'John', 'Red']

# Create a list from string using split()
groceries = "eggs, milk, cheese"
grocery_list = groceries.split(", ")
print(grocery_list)             # ['eggs', 'milk', 'cheese']
\end{lstlisting}
\end{examplebox}

\subsection{Mutability / 可变性（关键区别）}

\begin{concept}{Lists Are Mutable / 列表是可变的}
    Unlike tuples, lists are \textbf{mutable}. You can change the value at an index even after the list has been created. \\
    列表是\textbf{可变}的，你可以随时修改、添加或删除其中的元素。
\end{concept}

\begin{examplebox}{Mutability Demo / 可变性演示}
\begin{lstlisting}
grocery_list = ['eggs', 'milk', 'cheese']
grocery_list[1] = 'butter'
print(grocery_list)             # ['eggs', 'butter', 'cheese']
\end{lstlisting}
\end{examplebox}

\subsection{List Methods / 列表常用方法}

\begin{concept}{Why Methods? / 为什么需要方法？}
    Although you can add and remove elements with slice notation, list methods provide a more \textbf{natural and readable} way to mutate a list. \\
    虽然可以用切片增删元素，但列表方法更加直观易读。
\end{concept}

\begin{examplebox}{insert(), pop(), append(), extend() / 四大核心方法}
\begin{lstlisting}
mylist2 = []
mylist2.insert(0, 'A')          # Insert 'A' at index 0
print(mylist2)                  # ['A']

mylist2.insert(2, 'C')          # Insert at index 2 (beyond end → appended)
print(mylist2)                  # ['A', 'C']

mylist2.insert(1, 'B')          # Insert 'B' at index 1
print(mylist2)                  # ['A', 'B', 'C']

# pop(i): Remove and RETURN the value at index i
out = mylist2.pop(2)            # Removes 'C'
print(out)                      # 'C'
print(mylist2)                  # ['A', 'B']

# append(x): Add to the end
mylist2.append('D')
print(mylist2)                  # ['A', 'B', 'D']

# extend(iterable): Add multiple elements to the end
mylist2.extend(['E', 'F'])
print(mylist2)                  # ['A', 'B', 'D', 'E', 'F']
\end{lstlisting}
\end{examplebox}

\begin{warning}{pop() with Out-of-range Index / pop() 越界问题}
    Python raises an \texttt{IndexError} if you pass to \texttt{.pop()} an argument larger than the last index. \\
    如果传给 \texttt{pop()} 的索引超出范围，会抛出 \texttt{IndexError}。
\end{warning}

\subsection{Working with Lists of Numbers / 数字列表}

\begin{concept}{Built-in Functions for Numbers / 数字专用内置函数}
    Python provides convenient built-in functions for working with lists of numbers.
\end{concept}

\begin{examplebox}{sum(), min(), max() / 求和、最小值、最大值}
\begin{lstlisting}
mylist3 = [4, 17, 42, 5]

# Traditional way / 传统方式
total = 0
for number in mylist3:
    total = total + number
print(total)                    # 68

# Pythonic way / Python 风格
print(sum(mylist3))             # 68
print(min(mylist3))             # 4
print(max(mylist3))             # 42
\end{lstlisting}
\end{examplebox}

\subsection{The Reference Pitfall / 引用陷阱（重要考点）}

\begin{warning}{Memory Reference / 内存地址引用}
    Assigning one list to another (\texttt{list2 = list1}) does \textbf{NOT} create a copy. Both refer to the \textbf{same object in memory}! \\
    直接赋值只是起了个"别名"，修改其中一个，另一个也会跟着变。

    A variable name is really just a \textbf{reference} to a specific location in computer memory. Instead of copying all the contents, \texttt{large\_cats = animals} assigns the memory location referenced by \texttt{animals} to \texttt{large\_cats}. Both variables now refer to the \textbf{same object}.
\end{warning}

\begin{examplebox}{Right Way to Copy / 正确拷贝方式}
\begin{lstlisting}
animals = ["lion", "tiger"]

# WRONG way / 错误方式 — 只是别名
large_cats = animals
large_cats.append("Leopard")
print(large_cats)   # ['lion', 'tiger', 'Leopard']
print(animals)      # ['lion', 'tiger', 'Leopard'] — ALSO changed!

# RIGHT way / 正确方式 — 使用切片创建独立副本
animals = ["lion", "tiger"]
large_cats = animals[:]         # Creates a NEW list with same values
large_cats.append("Leopard")
print(large_cats)   # ['lion', 'tiger', 'Leopard']
print(animals)      # ['lion', 'tiger'] — unchanged!
\end{lstlisting}
\end{examplebox}

\subsection{Nesting Lists and Tuples / 嵌套结构}

\begin{concept}{Lists/Tuples Can Contain Lists/Tuples / 嵌套结构}
    Lists and tuples can contain lists and tuples as values. A \textbf{nested list} (or nested tuple) is a list or tuple that is contained as a value in another list or tuple. \\
    列表和元组可以互相嵌套。这在数据分析中常用于表示矩阵。
\end{concept}

\begin{examplebox}{Nested Structures / 嵌套示例}
\begin{lstlisting}
my_nested_list = [[1, 2], [3, 4]]
print(my_nested_list)           # [[1, 2], [3, 4]]
print(my_nested_list[1])        # [3, 4]

# Double index notation / 双重索引
print(my_nested_list[1][0])     # 3 (first element of second sublist)
\end{lstlisting}
\end{examplebox}

\subsection{Sorting Lists / 列表排序}

\begin{concept}{.sort() vs sorted() / 原地排序 vs 返回新序列}
    \begin{itemize}[itemsep=0.5em]
        \item \textbf{\texttt{list.sort()}}：Sorts the list \textbf{in-place} (modifies the original list). Returns \texttt{None}. Available for \textbf{lists only}.
        \item \textbf{\texttt{sorted(iterable)}}：Returns a \textbf{new sorted list} from any iterable. The original is unchanged. Works on \textbf{any iterable} (lists, tuples, strings, etc.).
    \end{itemize}
\end{concept}

\begin{examplebox}{Sorting Examples / 排序示例}
\begin{lstlisting}
# .sort() — in-place, list only
numbers = [1, 10, 5, 3]
numbers.sort()                  # Modifies 'numbers' directly
print(numbers)                  # [1, 3, 5, 10]

# sorted() — returns new list, works on any iterable
numbers = [1, 10, 5, 3]
new = sorted(numbers, reverse=True)
print(new)                      # [10, 5, 3, 1]
print(numbers)                  # [1, 10, 5, 3] — unchanged!

# Sort by key / 按自定义规则排序
colors = ["red", "yellow", "green", "blue"]
colors.sort(key=len)            # Sort by string length
print(colors)                   # ['red', 'blue', 'green', 'yellow']
\end{lstlisting}
\end{examplebox}

% ======================================================================
\section{Dictionaries / 字典 (Curly Braces: \texttt{\{k: v\}})}
% ======================================================================

\subsection{Key-Value Pairs / 键值对逻辑}

\begin{concept}{Mapping vs. Sequence / 映射 vs 序列}
    Dictionaries, like lists and tuples, store a collection of objects. However, instead of storing objects in a \textbf{sequence}, dictionaries hold information in pairs of data called \textbf{key-value pairs}. \\
    字典通过唯一的"键"（key）来存储和查找"值"（value），而不是靠位置索引。

    \begin{itemize}[itemsep=0.3em]
        \item \textbf{Key}：A unique name that identifies the value part of the pair. / 用于标识值的唯一名称。
        \item \textbf{Value}：The data associated with that key. / 与键关联的数据。
        \item Each key is separated from its value by a colon (\texttt{:}).
        \item Each key-value pair is separated by a comma (\texttt{,}).
        \item The entire dictionary is enclosed in curly braces \texttt{\{ \}}.
    \end{itemize}

    \textbf{Sequence types} (Lists/Tuples) use integers as indices. \textbf{Dictionaries} use unique keys as labels — fundamentally different!
\end{concept}

\subsection{Creating and Accessing Dictionaries / 创建与访问}

\begin{examplebox}{Dictionary Basics / 字典基础操作}
\begin{lstlisting}
# Creating a dictionary / 创建字典
capitals = {
    "France": "Paris",
    "Spain": "Madrid",
    "Italy": "Rome",
}
print(capitals)
# {'France': 'Paris', 'Spain': 'Madrid', 'Italy': 'Rome'}

# Empty dictionary / 空字典（两种方式）
dict_example = {}
dict_example = dict()

# Accessing values / 访问值
print(capitals['France'])       # 'Paris'
\end{lstlisting}
\end{examplebox}

\subsection{Adding and Removing Items / 增删条目}

\begin{concept}{Dictionaries Are Mutable / 字典是可变的}
    Like lists, dictionaries are mutable data structures. You can add and remove items dynamically. \\
    字典是可变的，可以动态增删条目。
\end{concept}

\begin{examplebox}{Adding and Deleting / 增删操作}
\begin{lstlisting}
capitals = {"France": "Paris", "Spain": "Madrid"}

# Adding a new item / 添加新条目
capitals['Germany'] = "Berlin"
print(capitals)
# {'France': 'Paris', 'Spain': 'Madrid', 'Germany': 'Berlin'}

# Deleting an item / 删除条目
del capitals["Spain"]
print(capitals)
# {'France': 'Paris', 'Germany': 'Berlin'}
\end{lstlisting}
\end{examplebox}

\subsection{Existence Check / 存在性检查}

\begin{warning}{KeyError When Key Doesn't Exist / 访问不存在的键会报错}
    Accessing a non-existent key raises a \textbf{KeyError}. Always use \texttt{in} to check first! \\
    访问不存在的键会抛出 \texttt{KeyError}，务必先用 \texttt{in} 检查！
\end{warning}

\begin{examplebox}{Safe Key Access / 安全的键访问}
\begin{lstlisting}
capitals = {"France": "Paris", "Spain": "Madrid"}

# Check before accessing / 先检查再访问
if "Spain" in capitals.keys():
    print(capitals["Spain"])    # 'Madrid'

# Simpler version / 更简洁的写法
if "Spain" in capitals:         # in 默认检查 keys
    print(capitals["Spain"])
\end{lstlisting}
\end{examplebox}

\subsection{Iterating Over Dictionaries / 字典遍历}

\begin{concept}{Different Ways to Iterate / 多种遍历方式}
    Dictionaries are iterable, but looping over them is different than looping over lists or tuples.
\end{concept}

\begin{examplebox}{Dictionary Iteration / 字典遍历方法}
\begin{lstlisting}
capitals = {"France": "Paris", "Spain": "Madrid", "Italy": "Rome"}

# Method 1: Loop over keys (default) / 默认遍历键
for key in capitals:
    print(key)
# Output: France, Spain, Italy (order may vary in Python < 3.7)

# Method 2: Access both key and value / 同时遍历键和值
for country in capitals:
    print(f"The capital of {country} is {capitals[country]}")

# Method 3: Using .items() — most "Pythonic" / 最 Python 的写法
for country, capital in capitals.items():
    print(f"The capital of {country} is {capital}")
# .items() returns (key, value) tuples that can be unpacked
\end{lstlisting}
\end{examplebox}

% ======================================================================
\section{Sets / 集合 (Curly Braces: \texttt{\{\}})}
% ======================================================================

\begin{concept}{What is a Set? / 什么是集合？}
    A set is a collection which is:
    \begin{itemize}[itemsep=0.3em]
        \item \textbf{Unordered / 无序}：Elements have no fixed position.
        \item \textbf{Mutable container / 容器可变}：You can add or remove elements after creation.
        \item \textbf{Hashable elements / 元素需可哈希}：Elements usually must be immutable values such as numbers, strings, or tuples.
        \item \textbf{Unindexed / 无索引}：You cannot access items by index (no \texttt{set[0]}).
        \item \textbf{No Duplicates / 不重复}：Sets cannot have two items with the same value. Duplicates are automatically removed!
    \end{itemize}
    集合是无序、无索引、不可重复的数据结构。集合本身可以增删元素，但元素需要是可哈希对象；写入相同值会被自动去重。
\end{concept}

\begin{examplebox}{Set Creation / 创建集合}
\begin{lstlisting}
myset = {"apple", "banana", "cherry"}
print(myset)                    # {'apple', 'banana', 'cherry'}
print(type(myset))              # <class 'set'>

# Duplicates are automatically removed! / 重复值自动去重
myset = {"apple", "banana", "cherry", "apple"}
print(myset)                    # {'apple', 'banana', 'cherry'} (no duplicate)
\end{lstlisting}
\end{examplebox}

\begin{warning}{Empty Set Requires set() / 空集合必须用 set()}
    \texttt{\{\}} creates an empty \textbf{dictionary}, NOT an empty set! \\
    \texttt{\{\}} 创建的是空\textbf{字典}，不是空集合！
    \begin{lstlisting}
empty_dict = {}                 # This is a dict!
empty_set = set()               # This is a set!
print(type(empty_set))          # <class 'set'>
    \end{lstlisting}
\end{warning}

% ======================================================================
\section{How to Pick a Data Structure? / 如何选择数据结构？}
% ======================================================================

\begin{table}[h]
\centering
\renewcommand{\arraystretch}{1.8}
\begin{tabular}{@{}p{4cm}p{4cm}p{4cm}@{}}
\toprule
\textbf{Use a List / 用列表} & \textbf{Use a Tuple / 用元组} & \textbf{Use a Dictionary / 用字典} \\ \midrule
Data has a natural order / 数据有自然顺序 & Data has a natural order / 数据有自然顺序 & Data is unordered, or order does not matter / 数据无序 \\
You need to update/alter data / 需要更新修改数据 & You will NOT update/alter data / 不需要修改数据 & You need to update/alter data / 需要更新修改数据 \\
Primary purpose is iteration / 主要用于遍历 & Primary purpose is iteration / 主要用于遍历 & Primary purpose is looking up values / 主要用于按键查找 \\ \bottomrule
\end{tabular}
\end{table}

% ======================================================================
\section{List Comprehensions / 列表推导式}
% ======================================================================

\begin{concept}{Short-hand for Loops / for 循环的优雅简写}
    A list comprehension is a succinct way to create a list from an existing iterable. It is a \textbf{short-hand for a for loop}. \\
    列表推导式是 \texttt{for} 循环的优雅简写形式，一行代码就能完成"遍历+转换+过滤"。
\end{concept}

\begin{examplebox}{Comprehension vs Traditional Loop / 推导式 vs 传统循环}
\begin{lstlisting}
numbers = (1, 2, 3, 4, 5)

# Traditional for loop / 传统 for 循环
squares = []
for num in numbers:
    squares.append(num**2)
print(squares)                  # [1, 4, 9, 16, 25]

# List comprehension / 列表推导式（一行搞定！）
squares = [num**2 for num in numbers]
print(squares)                  # [1, 4, 9, 16, 25]
\end{lstlisting}
\end{examplebox}

\begin{examplebox}{Comprehension with Type Conversion / 推导式 + 类型转换}
    列表推导式常用于将列表中的值转换为不同类型：
\begin{lstlisting}
# Convert strings to floats / 字符串列表 → 浮点数列表
str_numbers = ["1.5", "2.3", "5.25"]
float_numbers = [float(value) for value in str_numbers]
print(float_numbers)            # [1.5, 2.3, 5.25]
\end{lstlisting}
\end{examplebox}

\begin{concept}{Comprehension with Condition / 带条件的推导式}
    You can add an \texttt{if} condition at the end to \textbf{filter} elements:
\begin{lstlisting}
# Only square even numbers / 只对偶数求平方
even_squares = [num**2 for num in numbers if num % 2 == 0]
# Dictionary comprehension / 字典推导式
squares_dict = {num: num**2 for num in numbers}
    \end{lstlisting}
\end{concept}

% ======================================================================
\section{Quick Reference / 速查表}
% ======================================================================

\begin{table}[h]
\centering
\small
\renewcommand{\arraystretch}{1.3}
\begin{tabular}{@{}ll@{}}
\toprule
\textbf{Operation / 操作} & \textbf{Syntax / 语法} \\ \midrule
创建元组 & \texttt{t = (1, 2, 3)} \\
创建列表 & \texttt{lst = [1, 2, 3]} \\
创建字典 & \texttt{d = \{"a": 1\}} \\
创建集合 & \texttt{s = \{1, 2, 3\}} / \texttt{s = set()} \\
列表追加 & \texttt{lst.append(x)} \\
列表插入 & \texttt{lst.insert(i, x)} \\
列表弹出 & \texttt{lst.pop(i)} \\
列表扩展 & \texttt{lst.extend(iterable)} \\
列表排序(原地) & \texttt{lst.sort()} \\
列表排序(返回新) & \texttt{sorted(iterable)} \\
字典取值 & \texttt{d["key"]} / \texttt{d.get("key")} \\
字典添加 & \texttt{d["new\_key"] = value} \\
字典删除 & \texttt{del d["key"]} \\
检查键存在 & \texttt{"key" in d} \\
字典遍历 & \texttt{for k, v in d.items():} \\
列表推导式 & \texttt{[expr for var in iterable]} \\
带条件推导式 & \texttt{[expr for var in iterable if cond]} \\
print 分隔符 & \texttt{print(x, y, sep=",")} \\
print 结尾符 & \texttt{print(x, end="!")} \\ \bottomrule
\end{tabular}
\end{table}

\end{document}
